\documentclass[conference]{IEEEtran}
\IEEEoverridecommandlockouts
\usepackage{amsmath,amssymb,amsthm}
\usepackage{booktabs}
\usepackage{graphicx}
\usepackage{xcolor}
\usepackage{microtype}
\usepackage{url}
\usepackage[hidelinks]{hyperref}
\usepackage{cite}

\newcommand{\Q}{\mathbb{Q}}
\newcommand{\risk}[1]{L_D(#1)}
\newtheorem{proposition}{Proposition}
\newtheorem{theorem}{Theorem}
\newtheorem{assumption}{Assumption}

\begin{document}

\title{Detection Is Not Enforcement:\\
Evidence Repair and Risk-Gated Release in Tool-Using ML Pipelines}

\author{\IEEEauthorblockN{Michael Jerge}
\IEEEauthorblockA{\textit{Independent} \\ \url{https://github.com/mmjerge/galois-ml}}}

\maketitle

\begin{abstract}
Agentic ML pipelines act on evidence from tools they cannot verify.
Evaluating defenses against corrupted evidence is hard because the
correct outcome of a realistic agent task is rarely known, so harm is
scored by another model or a rubric. We construct a testbed in which
the correct outcome is a theorem: classification tasks from
arithmetic statistics (Galois groups of number fields from Dedekind
zeta coefficients) with a computable Bayes-optimal ceiling, an exact
per-input posterior, and a decidable predicate for which observations
are possible at all, so corruption harm is measured against an exactly
defined ceiling rather than a model's judgment.
Using it, we separate three questions that are usually conflated:
whether corrupted evidence is \emph{detectable}, whether its effects
can be \emph{removed}, and whether a pipeline's final artifact should
be \emph{released}. We show (i)~a deterministic verifier with stated
guarantee classes detects a decidable corruption class with zero false
positives, together with a restricted impossibility result for what no
label-symmetric verifier can detect; (ii)~in a sequential pilot with an identical detector and identical
detections, a warning-only policy was associated with a large
regression on a combined attack (median accuracy $0.90\to0.15$,
complete separation of runs) while silent verbatim repair of detected
responses was associated with near-clean performance; a prespecified,
randomized, interleaved replication is registered and under way;
(iii)~a fresh-holdout release gate bounds released risk relative to a
trusted fallback against \emph{arbitrary} upstream corruption,
including the provably undetectable class, with a trusted release
runtime that survived adversarial source audits. Guarantees are
independent of the number of agents and of agent disposition; they are
fallback-relative, conditional on stated trust assumptions, and do not
establish general semantic correctness.
\end{abstract}

\section{Introduction}

When a pipeline of language-model agents reaches a conclusion after
consulting tools, retrieved documents, or other agents, three
different things can go wrong with the evidence it acted on. The
evidence may be \emph{malformed}: values no honest source could
produce. It may be \emph{falsified but plausible}: real-looking values
attributed to the wrong things. Or the pipeline's own \emph{checks}
may be compromised, so that it cannot detect its errors. Defenses are
usually evaluated by asking whether a model \emph{notices} such
corruption, and the notice rate is reported as the security metric.

We argue, and then measure, that this metric answers the wrong
question. Detecting corruption is a property of a checker; whether the
detection changes the pipeline's output is a property of the
enforcement policy that acts on it; and whether the output is safe to
use is a property of a release decision that need not depend on
either. These are separate layers with separate guarantees, and a
pipeline can have the first without the second, or the third without
the first.

Measuring this requires knowing the right answer. We use a family of
synthetic tasks built from a theorem plus one stated modeling choice:
for a Galois number field of degree $d$, the Dedekind zeta coefficients
at an unramified prime are determined by the splitting type (a
theorem), and the Chebotarev density theorem fixes the asymptotic
frequency of each type over primes; our generator additionally draws
splitting types \emph{independently} across primes from that
distribution (the modeling choice). Under that fully specified law
each task has a Bayes ceiling that is exactly defined and computable
to Monte Carlo precision, a posterior defined on every
positive-probability input, and a decidable set of \emph{attainable}
coefficient values. The number theory is the instrument: it makes
correctness, support violation, and label inversion independently
checkable without a model-as-judge, and gives an exactly defined
ceiling against which harm is measured (the regrets we report are
empirical accuracies against a Monte Carlo estimate of that ceiling;
exact conditional regret is available per supported input but is not
the quantity in our tables). We do not claim the tasks resemble
deployed agent workloads, and we say explicitly what does and does not
transfer (\S\ref{sec:limits}).

\paragraph{Contributions}
\begin{enumerate}
\item \textbf{A controlled testbed for evidence corruption with
independently checkable outcomes} (\S\ref{sec:testbed}): tool-using
classification tasks with exact support constraints and posterior
risk under a stated synthetic law; corruption channels spanning
malformed outputs, recomputable falsehoods, and consistent
relabeling; and a restricted label-symmetry indistinguishability
result (Prop.~\ref{prop:imp}) for verifiers whose trusted observations
cannot distinguish class assignments. Validation against genuine
number fields (LMFDB) assesses mathematical fidelity, not
deployment-distribution equivalence.
\item \textbf{A policy-level evaluation separating detection from
mitigation} (\S\ref{sec:policy}): with an identical deterministic
detector and identical detection outcomes, two enforcement
policies---warn and silently repair---are associated in a sequential
pilot with outcomes ranging from a large regression to near-clean
performance. A prespecified, randomized, interleaved study of the
primary contrast is registered and in progress (\S\ref{sec:confirm});
pilot evidence is reported as exploratory association, not policy
causation.
\item \textbf{A layered enforcement implementation with distinct
guarantee classes} (\S\ref{sec:layers}): evidence checks, conditional
exact-response restoration, and a fresh-holdout release gate bounding
population excess risk relative to a trusted fallback
(Thm.~\ref{thm:gate}), with a trusted release runtime whose
obligations were adversarially audited. We expose the residual
accepted-corruption class with a defense-aware relabeling attack and
report both candidate and released-system performance.
\end{enumerate}

These are contributions about measurement, policy behavior, and the
composition of assurances. We do not claim new concentration bounds,
a defense without trusted resources, or provably trustworthy
multi-agent collaboration.

\section{Threat Model and Trusted Resources}\label{sec:threat}

\paragraph{System} One or more LLM agents, each with a fixed role
prompt, consult tools that expose a dataset: per-class coefficient
distributions, labeled examples, a list of structurally meaningful
indices, and a validation tool that executes a candidate rule on a
held split. Agents produce an executable classification rule
$h:\Q^N\to\{0,\dots,k-1\}$.

\paragraph{Adversary} Controls tool response channels: it may return
any bytes in place of an honest response, on any subset of calls, with
any per-call probability. Channels studied: \texttt{data}
(reattributed distributions, values attainable), \texttt{impossible}
(values the Euler-factor formula forbids), \texttt{verification}
(inflated validation, out-of-grammar), \texttt{verification\_stealth}
(inflated validation, grammar-conforming), \texttt{injection}
(appended operator-style instructions), \texttt{combined} (all,
mutually consistent), and \texttt{data\_adaptive} (consistent
relabeling across every evidence channel, validation left honest;
\S\ref{sec:adaptive}). A separate \texttt{saboteur} condition replaces
an agent's role prompt with a compromised one.

\paragraph{Trusted computing base} The mediator between tools and
agents; the dataset snapshot and any recomputation performed on it;
the certification sampler, labels, evaluator, and release mechanism;
a fixed fallback predictor. Out of scope: sandbox escape, credential
theft, resource exhaustion, and covert channels other than tool
content. These assumptions are substantive. Our contribution is to
separate the guarantees they support from guarantees that would
depend on agent behavior, not to remove them.

\section{Testbed}\label{sec:testbed}

For a Galois extension $K/\Q$ of degree $d$ and an unramified prime
$p$ with residue degree $f\mid d$ and $g=d/f$ primes above it,
\begin{equation}
a_{p^k}(K)=\begin{cases}\binom{g+k/f-1}{g-1}& f\mid k,\\ 0&\text{otherwise,}\end{cases}
\label{eq:euler}
\end{equation}
and $a_n$ is multiplicative. By Chebotarev, $f$ is distributed as the
order of a uniformly random group element. So the Galois group leaves
provable statistical fingerprints in the coefficient vector, and the
optimal decision rule is computable. Following Lee and
Lee~\cite{leelee2025}, we generate synthetic ``virtual fields'' from
this law for nine two- or three-class tasks (degree 4 through 25);
ceilings span $0.84$--$0.9998$ at $N{=}1000$. Security experiments use
sextic (C$_6$ vs.\ S$_3$, ceiling $0.9518$ at $N{=}1000$; $0.8991$ at
$N{=}200$) and degree-21 (C$_{21}$ vs.\ C$_7{:}$C$_3$, $0.9352$), which
leave headroom for harm to register.

\paragraph{The law, completely} Classes are uniform over the task's
group list. For each prime $p\le N$ ($N{=}1000$ in agent experiments,
$200$ in offline audits), a residue degree $f$ is drawn independently
from the class's Chebotarev measure (the distribution of element
orders); the visible signature $(a_{p^k})_{p^k\le N}$ follows
from~\eqref{eq:euler}; composite $a_n$ by multiplicativity. The
class-conditional likelihood of a row is the product over primes of
the pushforward mass of its signature (colliding signatures summed),
and the posterior is Bayes' rule with the uniform prior on rows of
positive mass; rows of zero mass have no posterior and are flagged
inadmissible. The Bayes ceiling is exactly defined by this law and
estimated by Monte Carlo ($2\times10^5$ Bernoulli samples of the Bayes
decision, binomial standard error $\approx5\times10^{-4}$ at a ceiling
of $0.95$); per-feature mutual information is computed in closed
form from the same pushforwards. The proven quartic criterion
classifies $200$ real LMFDB quartic fields ($100$ C$_4$, $100$
C$_2^2$, coefficients from true decomposition data) at accuracy
$1.000$; this validates the coefficient arithmetic, not the
independence assumption or the unramified restriction, which are
modeling choices. The held-out families (degrees 14, 21, 25) appear in
no publication we know of; this makes memorization of the \emph{task
instance} unlikely, not memorization of the underlying group theory.

\section{Three Layers, Three Guarantee Classes}\label{sec:layers}

\subsection{Evidence layer: deterministic verification}

A verifier interposed between tools and agents, importing nothing from
agent or attack code (standard library plus one LP call,
Prop.~\ref{prop:joint}). Two scopes must be kept apart: the
\emph{label-agnostic subverifier} (theorem and schema checks) knows
nothing that distinguishes candidate classes; the \emph{complete
mediator} additionally holds the trusted dataset snapshot and can
recompute the validation tool, which is class-anchored information.
Each check carries a guarantee class:

\begin{itemize}
\item \textbf{theorem}: \eqref{eq:euler} fixes the attainable values
of $a_{p^k}$ for \emph{any} Galois field of degree $d$ at unramified
primes, and multiplicativity gives the per-coordinate attainable set
of every composite $a_n$; a same-prime check requires one residue
degree to explain all visible powers of $p$ in a row; a
cross-coordinate check enforces $a_{mn}=a_ma_n$ on example rows where
every involved value renders losslessly. Values are checked as
membership in the renderer's \texttt{\%g}-image of the attainable set.
Sound for honest generator output; complete against reported values
outside the set. These are sound necessary checks, not a full
row-admissibility decision (the offline audit module has one).
\item \textbf{schema}: honest tools emit a fixed grammar; anything
outside it (appended instructions, label/class-name inconsistency,
forged refusals of valid requests) is flagged. Silent against
grammar-conforming forgeries.
\item \textbf{recomputation}: validation is deterministic under the
harness's controlled execution, so the mediator re-derives it; any
discrepancy is a proof of channel corruption. This is the only check
that catches the format-mimicking \texttt{verification\_stealth}
attack.
\end{itemize}

The theorem class also includes a \emph{joint} check on marginal
reports, which closes a gap the coordinatewise checks leave open:

\begin{proposition}[Same-prime closed-band realizability]\label{prop:joint}
Fix a prime $p$ and a class. For each visible power $k$, the honest
tool reports the class marginal of $a_{p^k}$, which is the image of
one distribution $q$ over residue degrees $f\mid d$ under
$f\mapsto R(b_{d,k}(f))$, where $b_{d,k}$ is \eqref{eq:euler} and $R$
the renderer. Call reported marginals $\{c_k\}$ \emph{closed-band
realizable} if some latent mixture $q$ has every rendered-value mass
within $\tau$ of the report. This holds iff the linear program
$q\ge0,\ \sum_f q_f=1,\ |\sum_{f:R(b_{d,k}(f))=v}q_f-c_k(v)|\le\tau$
for all $(k,v)$ is feasible (rendering collisions are summed in the
constraint rows).
\end{proposition}
``Realizable'' here means realizable by a \emph{population mixture}
whose rendered report lies within the closed rounding tolerance
(exactly $\pm\tfrac{1}{2000}$ for \texttt{\%.3f}); the finite-sample
version (integrality of $nq_f$ for a class of $n$ rows) is stronger
and not implemented. The LP has as many variables as divisors of $d$.
It is solved in floating point, but the verdict rests only on
certificates verified in exact rational arithmetic: a feasible point is
rationalized and every constraint re-checked, and a report is flagged
only on an exactly verified Farkas certificate of infeasibility
(obtained from a second LP); if neither verifies, or the solver
fails, the report is accepted. It is a sound infeasibility detector
with unresolved cases passing, not an exact decision procedure.
Duplicate rendered keys (distinct raw values sharing a \texttt{\%g}
image) are aggregated with the band widened accordingly. The check is
stateful across calls and transactional: a report is committed to the
store only after it passes, and under repair the committed report is
the honest replacement, so a rejected forgery cannot cause a later
honest report to be flagged. What a flag establishes is that the
accumulated evidence set is inconsistent with one sample---not, by
itself, that the current response is the corrupt one. It is class-agnostic---$q$ is free---so a consistent
relabeling passes, as Prop.~\ref{prop:imp} requires, while
individually attainable marginals that no mixture produces (e.g.\
$\Pr[a_p{=}d]=\tfrac12$ with $\Pr[a_{p^2}{=}\binom{d+1}{2}]=0$) are
refused.

Detection of the implemented fixtures is computed, not sampled
(Table~\ref{tab:matrix}), identical on the three tasks tested, and
scoped to those fixtures: it is not completeness for everything a
channel could carry. Soundness (``honest output is never flagged'')
is proved for the arithmetic checks (Lean, below) and tested against
the renderer for the grammar checks. The label-agnostic subverifier
is kept free of class-distinguishing knowledge because a subverifier
with it would be an oracle for the task and would void the
measurement.

\begin{table}[t]
\centering\small
\caption{Deterministic detection of the implemented attack fixtures
(computed; identical on all tasks). Honest output is never flagged.}
\label{tab:matrix}
\begin{tabular}{lcl}
\toprule
channel & detected & guarantee \\
\midrule
none (control) & no & --- (soundness) \\
data / data\_adaptive & \textbf{no, by design} & --- (Prop.~\ref{prop:imp}) \\
verification & yes & schema \\
verification\_stealth & yes & recomputation \\
injection & yes & schema \\
combined & yes & schema \\
impossible & yes & theorem \\
\bottomrule
\end{tabular}
\end{table}

\begin{proposition}[Restricted label-permutation indistinguishability]
\label{prop:imp}
Let a verifier observe class-indexed evidence transcripts, with
trusted knowledge invariant under a permutation $\pi$ of class labels,
and let the attack consistently apply $\pi$ to every class attribution
in every observed transcript, changing nothing else. If the verifier
has false-positive probability at most $\alpha$ in every admissible
honest world, its detection probability against the permuted
transcript is at most $\alpha$.
\end{proposition}
\begin{proof}
The attacked transcript has exactly the distribution of an honest
transcript in the world $\pi w$, which is admissible by
label-invariance of the trusted knowledge; apply the false-positive
bound in that world.
\end{proof}

The proposition is a boundary, not a defense. It does \emph{not}
apply when the verifier has trusted labeled examples, exact labeled
recomputation, or observes partial relabelings inconsistent across
channels. Its detection matrix row is pinned by a regression test that
fails if the verifier ever acquires class-distinguishing knowledge.

\subsection{Protocol layer: enforcement policies}

Given a detection, the mediator can \emph{flag} (return the response
verbatim with a diagnostic appended), \emph{flag with routing}
(diagnostic names the failed source and the channels that remain
trustworthy), \emph{block} (return a fixed code), or \emph{repair}
(return the trusted recomputation verbatim, with no notice). Block
gives behavior-independent \emph{containment} of the corrupt bytes;
only repair admits a behavior-independent \emph{restoration}
guarantee:

\begin{proposition}[Exact mediated-execution coupling]\label{prop:repair}
Consider a channel-attacked and a channel-clean execution with the
\emph{same} agent programs, initial prompts, orchestration policy, and
exogenous inputs (so any compromised agents are present in both), with
randomness modeled by coupled tapes. Let $\mathcal R$ relate their
trusted environment states, initially satisfied, such that whenever
agent states and visible histories agree and environment states are
$\mathcal R$-related: (1)~the next observable agent action is the
same; (2)~the attacked execution's mediated response to it equals the
clean execution's response; (3)~environment states remain
$\mathcal R$-related afterwards; (4)~termination and artifact
extraction agree. If no other attacker-controlled observation or state
transition reaches the agents, then every finite visible transcript
prefix is identical, termination occurs at the same step, and final
artifacts are identical.
\end{proposition}
\begin{proof}
Induction on steps; (1)--(3) preserve the invariant, (4) closes.
\end{proof}
Detect-and-repair establishes (2) provided recomputation uses the
correct snapshot and request and has no divergent side effects; for
stateless deterministic tools on a fixed snapshot, $\mathcal R$ is
equality. Two limits are part of the statement: the reference is the
\emph{channel-clean} execution with the same agents, so repair removes
covered channel corruption even for badly behaved agents but does not
turn a compromised agent into an honest one; and the conclusion is
equality to that reference, not correctness or low risk.

The proposition is conditional on complete detection and exact
restoration; tests verify byte equality of attacked-repaired and honest
responses on every implemented detected fixture, which establishes
hypothesis (2) for those fixtures and not for arbitrary bytes. The
\texttt{data} channels are undetected and therefore unrepaired; that
residual is what the third layer covers.

\paragraph{The three guarantees, kept separate} The layers compose
into three statements with different characters, which we deliberately
do not merge into one bound. (i)~\emph{Deterministic restoration}
(Prop.~\ref{prop:repair}): for the covered channel class, the attacked
execution equals the channel-clean reference with the same agents;
clean-relative, not a risk statement. (ii)~\emph{Statistical release
safety} (Thm.~\ref{thm:gate}): conditional on all pre-certification
randomness, for arbitrary upstream corruption including all $N$ agents
compromised and attacks outside the covered class,
$\Pr[\risk{h_{\mathrm{out}}}>\risk{h_0}+\epsilon]\le\delta$;
fallback-relative, per certification, with $K$ the only
proposal-count dependence, and no bound on
$\Pr[\text{unsafe}\mid\text{a candidate was released}]$. (iii)~\emph{Label-symmetry
boundary} (Prop.~\ref{prop:imp}): for the label-agnostic subverifier
only; a limit on detection, not on harm, and not on the gate's ability
to refuse the resulting proposal.

\paragraph{What is machine-checked} Four statements are formalized in
Lean~4 without Mathlib (\texttt{lean/}, 21 theorems, no \texttt{sorry}),
each mapped to a claim: (i) the attainable set of~\eqref{eq:euler}
over divisors of $d$, with the theorem that $0$ is attainable iff some
divisor of $d$ does not divide $k$ (the arithmetic support lemma
\emph{conditional on} the Euler formula---the number theory itself is
not formalized; this is the fact our first implementation violated);
(ii) Prop.~\ref{prop:imp} in its zero-false-positive form, as an
abstract statement about a verifier, honest set, and relabeling
action; (iii) Prop.~\ref{prop:repair} as a trace induction over a
deterministic transition system; (iv) the deterministic clause of
Theorem~\ref{thm:gate}: on the concentration event, any valid release
has risk $\le\risk{h_0}+\epsilon$, including that $\epsilon\ge0$ is
necessary and that fail-closed is always safe. The concentration
inequality is a hypothesis in (iv). None of these connects to the
parser or runtime code; they fix the mathematics the code is tested
against. The \texttt{data} channels are
undetected and therefore unrepaired; that residual is what the third
layer covers.

\subsection{Output layer: fallback-relative release}

The agent's artifact is executable. Its risk on fresh data is
checkable regardless of how the agent reasoned or how corrupt its
evidence was.

\begin{theorem}[Fresh-holdout gate]\label{thm:gate}
Let $\mathcal F$ contain all pre-certification randomness. Let
$h_1,\dots,h_K$ ($K\ge1$) be candidate predictors and $h_0$ a fallback,
all $\mathcal F$-measurable, pointwise, with identical certification
and deployment semantics, and let $H=((X_i,Y_i))_{i=1}^m$ satisfy
$H\mid\mathcal F\sim D^m$, $m\ge2$, with 0--1 loss $\ell$ and
$0<\delta<1$. Let $Z_{ji}=\ell(h_j;X_i,Y_i)-\ell(h_0;X_i,Y_i)\in\{-1,0,1\}$,
$\hat\Delta_j=\frac1m\sum_i Z_{ji}$, $V_j$ the unbiased sample
variance of $(Z_{ji})_i$, and
\begin{align}
r^{H}&=\sqrt{2\log(2K/\delta)/m},\label{eq:hoeff}\\
r^{B}_j&=\sqrt{2V_j\log(4K/\delta)/m}+\tfrac{14\log(4K/\delta)}{3(m-1)},\label{eq:bern}
\end{align}
$r_j=\min(r^H,r^B_j)$. A release rule may depend arbitrarily on $H$ but
releases only $h_0$ or a candidate with $\hat\Delta_j+r_j\le\epsilon$,
$0\le\epsilon\le1$. Then, conditional on $\mathcal F$ and hence
unconditionally,
$\Pr[\risk{h_{\mathrm{out}}}>\risk{h_0}+\epsilon]\le\delta$, and
$\mathbb{E}[(\risk{h_{\mathrm{out}}}-\risk{h_0})_+]\le\epsilon+\delta(1-\epsilon)$.
\end{theorem}
\begin{proof}
Work conditionally on $\mathcal F$. Spend $\delta/2$ on each inequality.
One-sided Hoeffding on $[-1,1]$ with a union bound over $K$ gives
$\risk{h_j}-\risk{h_0}\le\hat\Delta_j+r^H$ for all $j$ w.p.\
$\ge1-\delta/2$; the empirical-Bernstein inequality of Maurer and
Pontil~\cite{maurerpontil2009} (range $2$) gives the same with $r^B_j$
w.p.\ $\ge1-\delta/2$. On the intersection $E$, of probability
$\ge1-\delta$, every accepted $j$ has
$\risk{h_j}-\risk{h_0}\le\hat\Delta_j+r_j\le\epsilon$, and rejection
releases $h_0$ with excess $0\le\epsilon$; this is the first claim. For
the second, the positive excess is at most $\epsilon$ on $E$ and at
most $1$ on $E^c$, and $\Pr(E^c)\le\delta$, so its expectation is at
most $\epsilon+\delta(1-\epsilon)$. Average over $\mathcal F$.
\end{proof}
Repeated certifications each carry their own $\delta$; a campaign of
$T$ releases has failure probability at most $\sum_t\delta_t$, and
re-certifying until something passes is not covered.

The variance term is what makes the gate usable against a strong
fallback. Paired 0--1 differences are mostly zero when the two
predictors are close, so $V_j\ll1$ and $r^B\ll r^H$. A worked
instance, with the paired counts that determine everything: a genuine
agent rule (sextic, clean pilot run) against the theorem fallback on
one fresh $m{=}2000$ holdout gives $(n_+,n_0,n_-)=(72,1906,22)$, i.e.\
disagreement on $4.7\%$ of rows, $\hat\Delta=0.025$, $V=0.046$
(holdout accuracies $0.9345$ vs.\ $0.9595$). With $K{=}1$,
$\delta{=}0.05$: $r^H=\sqrt{2\ln40/2000}=0.0607$ and
$r^B=\sqrt{2\cdot0.046\ln80/2000}+14\ln80/(3\cdot1999)=0.0245$, a
$2.5\times$ reduction, and $\hat\Delta+r^B=0.0495\le\epsilon=0.05$: the
rule certifies, narrowly. Under a Hoeffding-only gate it is refused; a
range-only \emph{sufficient} sample size for $95\%$ acceptance of a
candidate with excess $\Delta=0.031$ is
$m\ge[\sqrt{2\ln(2K/\delta)}+\sqrt{2\ln 20}]^2/(\epsilon-\Delta)^2
\approx7.4\times10^4$; the actual power under the ternary law is
higher and this is not a necessity claim. Every gate number in this
paper is from the rule as stated (min of the two radii, $K$ given per
experiment); \S\ref{sec:gateoc} reports both gates side by side.

The bound has no factor in the number of agents $N$ or the number of
compromised agents $k\le N$: all agents may be compromised and submit
arbitrary code. It is \emph{fallback-relative}: an adversary can force
fallback and destroy the agents' improvement over $h_0$ at zero risk
cost, which is an availability loss, not a safety loss. Bounding
degradation relative to a \emph{clean} agent needs a bound on
$\risk{h_0}-\risk{h_{\mathrm{clean}}}$, which the gate does not supply.
The theorem assumes a fixed pointwise predictor; a batch program with
cross-row state is outside it.

\paragraph{Release runtime} A single-use lifecycle
(\textsc{open}$\to$\textsc{frozen}$\to$\textsc{certifying}$\to$
\textsc{certified}$\to$\textsc{released}, \textsc{failed} terminal):
candidates hashed before certification randomness exists; holdout seed
drawn from OS entropy after freeze, held privately, disclosed only in
the audit record; candidates executed in a spawned process receiving
features only, replying over a bounded non-pickle byte protocol;
label-alphabet validation of candidate and fallback; immutable private
authorization records; released candidates executing through the same
isolated path as certification. Two adversarial source audits found
and we fixed: mutable returned decisions controlling release,
exception-unsafe state, pickle deserialization from the untrusted
worker, missing label validation, and a divergent deployment path.
Remaining obligations for the surrounding system (executor manifest,
resource limits, protected deployment sink, signed audit storage) are
stated, not claimed (\S\ref{sec:limits}).

\section{Policy Results}\label{sec:policy}

\subsection{Pilot (exploratory)}

Claude Sonnet~4.5 via Amazon Bedrock; sextic; theory withheld from
the brief so the tool channel is load-bearing; temperature $0.7$;
$n=5$ per condition; arms run sequentially. Median accuracy
(Table~\ref{tab:threearm}).

\begin{table}[t]
\centering\small
\caption{Pilot: same detector, three policies. Median accuracy,
$n=5$/cell except flag/\texttt{data} ($n=4$; one API failure,
retained in the supplement), sextic (ceiling $0.9518$). Exact
one-sided permutation $p$ on accuracy sums, unpaired; nominal and
exploratory---sequential arms do not establish exchangeability.}
\label{tab:threearm}
\begin{tabular}{lccccc}
\toprule
channel & unguarded & flag & repair & $p$(R$>$F) & $p$(R$>$U)\\
\midrule
none & 0.913 & 0.925 & 0.920 & --- & ---\\
data & 0.907 & 0.898 & 0.907 & 0.48 & 0.62\\
injection & 0.925 & 0.920 & 0.925 & 0.50 & 0.84\\
impossible & 0.920 & 0.728 & 0.925 & 0.020 & 0.14\\
combined & 0.898 & \textbf{0.153} & \textbf{0.910} & \textbf{0.004} & 0.028\\
verif.\_stealth & --- & 0.890 & 0.925 & 0.032 & ---\\
\bottomrule
\end{tabular}
\end{table}

Under \texttt{flag}, every poisoned response on a detected channel was
flagged (deterministic; per-run counts in the repository), yet no
channel with an unguarded control had a higher guarded median, and
\texttt{combined} collapsed: in all five runs the final conjecture
states the correct $f{=}2$ criterion attached to the wrong class
(coded from the returned JSON conjecture and rule; the transcripts are
released). Flagging retains the original bytes at the response level,
so guarded agents received the original evidence plus a correct
diagnostic; transcripts show them abandoning the flagged channel for
surface heuristics. Under \texttt{repair}, every detected channel's
median lay within the clean band. Complete separation of the ten run
accuracies on \texttt{combined} repair-vs-flag ($p=1/\binom{10}{5}$)
is the contrast that motivates the confirmatory study; the others are
unresolved at this $n$.

What this does \emph{not} show: that warnings \emph{cause} harm (arms
were sequential, and the flag-vs-unguarded question is a separate
contrast), that warnings fail in general (one message design, one
model), that repair equals clean (a nonsignificant difference is not
equivalence), or that inverted rules defeat review in general (five
transcripts).

\subsection{Prespecified confirmatory study}\label{sec:confirm}

Protocol registered before data collection (git tag
\texttt{confirmatory-protocol-v1}): $3$ policies $\{$unguarded, flag,
repair$\}$ $\times$ $\{$\texttt{combined}, \texttt{impossible},
\texttt{none}$\}$ $\times$ $\{$sextic, degree-21$\}$ $=18$ cells,
$n=10$ runs per cell, all $180$ runs randomly interleaved from a
committed seed so policy is not confounded with time. Both endpoints
per run: raw pipeline accuracy, and the accuracy of the system a gated
deployment would have \emph{released} (certified through the release
runtime against a majority fallback, $K{=}1$, $m=20{,}000$,
$\epsilon=\delta=0.05$, audited on an independent $20{,}000$-sample
draw). Primary: sextic/\texttt{combined}, repair vs.\ flag, raw
endpoint, one-sided exact permutation, $\alpha=0.05$. Registered
secondaries, Holm-corrected as one family of four: sextic
\texttt{combined} repair vs.\ unguarded; sextic \texttt{impossible}
repair vs.\ flag; degree-21 \texttt{combined} repair vs.\ flag; the
released-endpoint version of the primary. Invalid runs (API failures)
are retried once and otherwise reported missing; clean channels are
reported descriptively as a drift check.

Two questions are \emph{not} in the registration and will be reported
as exploratory: whether flag is worse than unguarded (which is what
``warnings harmed agents'' would require), and the competent-advisory
baseline. For the latter we have implemented \texttt{flag\_routed}
(names the failed source and the channel that remains independently
verified, under identical trust assumptions and with no task answers);
it is not an arm of the registered study and will be run as a
separately registered follow-up. If routed warnings recover
performance, the contribution sharpens into a distinction between
behavior-dependent mitigation and the conditional restoration
guarantee of Prop.~\ref{prop:repair}; it does not weaken the
architectural claim, which never rested on warnings failing. Any
amendment to the protocol after data collection began will be
reported as a timestamped amendment.

\section{Defense-Aware Attacker and the Residual Class}\label{sec:adaptive}

Against repair, an attacker who adds a detectable marker causes the
defender to restore the very evidence it wanted to falsify. The
defense-aware play is one consistent label permutation across every
evidence channel, exact grammar, attainable values, class names
matched---and honest validation, since a validation lie is caught by
recomputation. This channel (\texttt{data\_adaptive}) is a
defense-aware attacker, not a proven optimum over the accepted-output
space.

End to end, offline (sextic, $N{=}200$, ceiling $0.8991$): the
label-agnostic subverifier, with every check enabled, detects $0$ of
the attack's responses---consistent with Prop.~\ref{prop:imp}, whose
hypotheses that subverifier satisfies. The complete mediator's
recomputation check is class-anchored and so is \emph{not} covered by
the proposition; it detects nothing here for a different reason: the
attack leaves validation honest, so there is no discrepancy to
recompute. A credulous learner (decision tree fit to the permuted
evidence) lands at true accuracy $0.240$ on $600$ test rows
($1-$ceiling $=0.101$); the release gate ($K{=}1$, $m{=}2000$,
theorem fallback) refuses it, so the released system is the fallback
and its excess on an independent $20{,}000$-sample audit is $0$ by
construction---evidence that the gate refused a catastrophic rule, not
of population-risk control.

Live agents (sextic, $N{=}1000$, $n=5$ per arm, $600$ test rows per
run, medians reported): clean $0.895$ (range $0.887$--$0.925$),
attacked-unguarded $0.920$ ($0.867$--$0.945$), attacked-repair $0.890$
($0.867$--$0.925$); $0$ detections by the guard in the $5$ guarded
runs across $142$ poisoned responses. The attack evaded the checks but
produced no clear degradation. Transcripts are consistent with two orientation anchors
the attacker cannot permute---trusted validation feedback (agents
iterate candidate rules against the honest validation tool) and
pretrained knowledge of the named groups' element orders---and this
experiment does not separate them. The structural statement we can
make: under trusted recomputation and complete mediation, a false
validation result cannot persist at the agent interface; whether that
\emph{disarms} the attack depends on agent behavior this study did
not isolate. The separating experiment is a $2\times2$ factorial
under the same attack---$\{$named groups, opaque per-run
labels$\}\times\{$validation available, unavailable$\}$, with final
class orientation as a prespecified mechanism endpoint---which we have
designed and not yet run. Our reading of the transcripts is that the
validation anchor is doing most of the work (agents revise after
failed validations and rarely cite group orders before testing), but
we hold that as a hypothesis.

The residual accepted-corruption class is larger than this fixture:
altered attainable frequencies, invented examples, grammar-conforming
error text, anything the parser accepts unchanged. The gate's bound is
what covers that entire class, fallback-relative.

\section{Gate Operating Characteristics}\label{sec:gateoc}

All numbers here use the gate of Theorem~\ref{thm:gate} as stated
($r=\min(r^H,r^B)$, $\delta/2$ each), with $K{=}1$ per certification
(each artifact certified separately on its own fresh holdout),
$\epsilon=\delta=0.05$. An earlier draft mixed these with a
Hoeffding-only gate; those numbers are superseded.

\paragraph{Retrospective integration study} The 89 archived pilot
rules were certified through the release runtime, $m=2000$, against
two fallbacks. Majority fallback (risk $\approx0.5$): $82/89$
released; the seven refusals were every rule at or below chance (six
inversions at $0.10$--$0.16$ holdout accuracy, one at exactly
$0.500$). Theorem fallback (risk $\approx0.049$): $8/89$ released
(five clean-channel rules, two \texttt{injection}, one
\texttt{verification\_stealth}), with a maximum released holdout
excess of $+0.026<\epsilon$; under a Hoeffding-only gate the same
artifacts gave $0/89$. That released artifacts satisfy the
certification inequality on the certification sample is an algebraic
property of the rule, not evidence about population risk. The
theorem-fallback regime is a safety-floor experiment: agents at
$\approx0.92$ against a $0.951$ fallback rarely clear
$\hat\Delta+r\le\epsilon$ at $m=2000$, and their doing so is not
evidence that they add value over the fallback.

\paragraph{Operating characteristics with independent audit} Boundary
candidates are the fallback with predictions flipped on rows selected
by a deterministic per-row hash at rate $q$; their population excess
$\Delta$ is \emph{estimated} on an independent $10^5$-row sample
(standard error $\le0.002$), not computed exactly. Thirty fresh
$m$-row holdouts per cell; every released artifact's excess is judged
against that $\Delta$ estimate. Results:

\begin{center}\small
\begin{tabular}{rccc}
\toprule
$\Delta$ (est.) & accept, $m{=}2000$ & accept, $m{=}20{,}000$ & viol./released \\
\midrule
$0.000$ & $30/30$ & $30/30$ & $0/60$ \\
$0.016$ & $30/30$ & $30/30$ & $0/60$ \\
$0.032$ & $6/30$ & $30/30$ & $0/36$ \\
$0.044$ & $0/30$ & $19/30$ & $0/19$ \\
$0.055$ ($>\epsilon$) & $0/30$ & $0/30$ & --- \\
$0.079$, $0.159$ & $0/30$ & $0/30$ & --- \\
\bottomrule
\end{tabular}
\end{center}

With 30 replicates a $30/30$ cell has a $95\%$ Clopper--Pearson lower
bound of $0.88$ and a $0/30$ cell an upper bound of $0.12$. In the observed cells, acceptance is complete well inside the budget,
degrades as $\Delta$ approaches the effective threshold $\epsilon-r$,
and is zero in every cell with $\Delta>\epsilon$ (no cell sits exactly
at $\epsilon$); no released artifact exceeded $\epsilon$ excess by the
independent estimate. These are per-certification results at $K{=}1$;
they do not bound the probability that \emph{some} certification in
the campaign was unsafe, which is $\le\sum\delta_t$. The $\Delta{=}0$ row (candidate identical to the
fallback, zero disagreement variance) certifies at $m=2000$, as the
Bernstein radius requires and a Hoeffding-only gate would not.

\section{Related Work}
\paragraph{Reference monitors and complete mediation} The mediator of
\S\ref{sec:layers} is a reference monitor in Anderson's
sense~\cite{anderson1972}: every tool response passes through one
trusted point, and Prop.~\ref{prop:repair}'s complete-mediation
hypothesis is Saltzer and Schroeder's principle~\cite{saltzer1975}
applied to evidence rather than to access. What is new is not the
monitor but what it checks (arithmetic admissibility rather than
authorization) and what it does on failure (restoration rather than
denial). \paragraph{Prompt injection and tool-output defenses} Indirect
prompt injection through retrieved or tool-returned content is
established~\cite{greshake2023}. Design-level defenses such as
CaMeL~\cite{camel2025} separate control flow from untrusted data with
capability tracking so that injected content cannot acquire authority;
compromised-agent architectures (SAGA-BFT, MUZZLE, ACE~\cite{nds2})
work at the permission and containment boundary. These systems
principally constrain the \emph{authority} of untrusted content; our
checks constrain specified arithmetic properties of tool evidence,
which are necessary (not sufficient) for its truth. The protections
are complementary: capability confinement alone does not establish
those properties. \paragraph{Baseline-relative certification}
Releasing a candidate only when a fresh sample certifies it against a
baseline is the structure of high-confidence and safe policy
improvement~\cite{thomas2015,laroche2019}, of selective
classification~\cite{geifman2017}, and of risk-controlling
prediction~\cite{bates2021,angelopoulos2021}. Theorem~\ref{thm:gate}
is an instance of that pattern with paired 0--1 losses and an
empirical-Bernstein radius~\cite{maurerpontil2009}; we claim no new
inequality. Our contribution is placing the gate at the exit of an
adversarial pipeline whose upstream may be entirely compromised, and
showing which corruption classes it is the \emph{only} layer that
covers. \paragraph{Behavioral studies of multi-agent systems} Anthropic's
\emph{Patterns and problems in multiagent systems}~\cite{anthropic2026}
provides observational context through studies of trust, conformity,
and coordination. Its lie-detection evaluation and ours share a
comparison pattern: agent behavior is assessed against a credulous
policy and a ground-truth benchmark. The tasks, information
structures, and benchmark semantics differ; our setting additionally
permits analytic ground truth and formal bounds. The report concludes
that models may understand source incentives yet lack ``a disposition
to act on that knowledge without prompting.'' We take this as
motivation for a narrower architectural question: which assurances can
be enforced without relying on upstream agents' judgment? Our repair
mechanism and release gate enforce distinct evidence and release
conditions under explicit trust and sampling assumptions. They do not
establish improved coordination, incentive alignment, or information
pooling. Its conformity observations are qualitatively consistent with
our earlier measurement, on this testbed, that the error sets of five
independently run agents overlapped $14$--$26\times$ more than the
Jaccard index implied by independent errors at their rates
($pq/(p+q-pq)$); the two measure different phenomena, and that
measurement motivates but is not a premise of the agent-count-independent
guarantees.

\paragraph{Discussion: enforcement rather than discretionary response}
Our response to the disposition gap is deliberately limited. Where a
tool-evidence property admits a tractable, sound check from trusted
inputs, the corresponding assurance need not depend on an agent
noticing a warning or choosing to honor it. The release gate provides a
separate, statistical guarantee relative to a baseline; it does not
guarantee detection of every harmful candidate. Disposition
independence therefore concerns enforcement at the trusted boundaries,
not the behavior or integrity of the entire system. The report's
observation of price matching through a public listings board after
direct communication was removed illustrates a related limitation:
task outputs can remain a medium for coordination. Our formal boundary
applies to the specified verifier and its admissible or
observationally indistinguishable output class; we do not establish
that the report's experiment satisfies those hypotheses, nor claim that
content verification prevents collusion. Likewise, repairing tool
evidence does not itself ensure that agents elicit, transmit, or use
pivotal private information. Lee and Lee~\cite{leelee2025} established the arithmetic
criteria we use as ground truth. Selective classification and
fallback-relative certification are standard; our contribution is
their composition with mediation and repair, and the separation of
which assurance each supplies.

\section{Limitations}\label{sec:limits}
\emph{Scope of exactness.} Theorem-class checks are exact for the
unramified synthetic law; ramified primes of real fields fall outside
them. Deployment distributions are out of scope. \emph{What
transfers.} The logical structure---complete mediation, trusted
restoration, frozen proposals, fresh-holdout certification---and not
the observed rates of warning-induced failure, attack effectiveness,
or the cheap availability of IID labels and exact recomputation.
\emph{Trusted resources.} Recomputation and labeled certification data
are substantive; we compare policies under the same budget and report
acceptance and fallback frequency alongside risk. \emph{Runtime.}
Source hashing is not semantic artifact identity; pointwise semantics
are assumed; resource limits, a protected deployment sink, and signed
audit storage are the surrounding system's obligations.
\emph{Empirical.} One model family; $n=5$ sequential pilot arms with
the confirmatory study in progress; the anchor factorial and the
routed-warning baseline are designed but not run. No pilot contrast
is causal; the \texttt{combined} repair-vs-flag separation is the one
we regard as strong enough to have motivated a registered
replication. \emph{Related work.} The individual principles---complete
mediation, trusted replacement, baseline-relative certification---are
established; our claim is their composition and instrumentation, and
the related-work discussion here is briefer than the final version
must be.

\section{Conclusion}
A checker is not a defense until something binding acts on its
verdicts, and the thing that acts should not be the agent's
disposition. In this testbed, the same detector with the same
detections was associated with a large regression under one
enforcement policy and near-clean performance under another, and a
relabeling attack the label-agnostic checks cannot see was refused by
a release gate that never looked at the evidence at all. The
guarantees are conditional and fallback-relative, they do not depend
on how many agents there are or on what any of them believes, and the
behavioral finding awaits its registered replication.

\appendix
\section{Claim Ledger and Experiment Status}

\begin{table*}[t]
\centering\small
\caption{Claim ledger: what is claimed, on what assumptions, with what
evidence, and its current status.}
\label{tab:ledger}
\begin{tabular}{p{3.6cm}p{4.2cm}p{4.6cm}p{3.4cm}}
\toprule
claim & assumptions & evidence & status \\
\midrule
Attainable-value check is sound and complete for reported values outside the set &
unramified synthetic law; renderer image & Lean theorem (\texttt{Attainable.lean}); soundness sweep on all 9 tasks & proved (arithmetic), tested (renderer) \\
Same-prime joint realizability iff LP feasible & population mixture; closed rounding band & exact rational certificates; 3 tests & implemented; finite-sample version not \\
Label-agnostic subverifier cannot detect consistent relabeling (Prop.~\ref{prop:imp}) &
closure of honest set under $\pi$; zero FP & Lean (\texttt{Indistinguishability.lean}); regression test; 0/142 live &
proved (zero-FP form); $\alpha$-form on paper only \\
Exact repair $\Rightarrow$ clean transcript (Prop.~\ref{prop:repair}) &
complete mediation; exact recomputation; coupled randomness; no side channels &
Lean (\texttt{Repair.lean}); byte-equality tests on all detected fixtures &
proved conditional; hypotheses tested on fixtures \\
Gate: $\Pr[\text{excess}>\epsilon]\le\delta$ (Thm.~\ref{thm:gate}) &
frozen pointwise predictors; fresh IID holdout; trusted labels/evaluator & standard inequalities; Lean deterministic clause; OC study, 0/175 released violations & proved; fallback-relative; pointwise semantics assumed \\
Release runtime enforces the gate & single trusted process & 17 tests incl.\ mutation attack, lifecycle, IPC, timeout & implemented; executor manifest, resource limits, sink, signed audit pending \\
Warning-only policy associated with regression on \texttt{combined} & one model, one message & pilot $n{=}5$, complete separation; orientation coding 5/5 inverted & exploratory; replication registered, running \\
Repair associated with near-clean performance & same & pilot; repair vs flag $p=1/252$ & exploratory; not an equivalence claim \\
Warnings \emph{harm} agents (flag $<$ unguarded) & --- & not registered & exploratory only \\
Competent warning (\texttt{flag\_routed}) & --- & implemented, not run & pending \\
Adaptive attack harmless to live agents; mechanism & --- & $n{=}5$/arm; two candidate anchors & descriptive; factorial designed \\
\bottomrule
\end{tabular}
\end{table*}

Table~\ref{tab:ledger} consolidates every claim with its assumptions,
evidence, and status. The supplementary artifact
(\texttt{paper/satml/supplement/}) contains run-level outcomes for all
135 pilot and adaptive-attack runs (\texttt{pilot\_runs.csv}), the
mechanical orientation coding behind the inversion finding
(\texttt{orientation\_coding.csv}: each returned rule is executed on
fresh synthetic data and labeled correct/inverted/neither by whether
it scores $>0.75$ against true or against permuted labels---flag/%
\texttt{combined} $5/5$ inverted, repair/\texttt{combined} $5/5$
correct, unguarded/\texttt{combined} $3$ correct, $1$ inverted, $1$
neither), and the per-run gate refusal manifest
(\texttt{gate\_refusals.csv}).

\begin{thebibliography}{19}
\bibitem{anderson1972} J. P. Anderson, ``Computer security technology
planning study,'' ESD-TR-73-51, USAF Electronic Systems Division, 1972.
\bibitem{saltzer1975} J. H. Saltzer and M. D. Schroeder, ``The
protection of information in computer systems,'' \emph{Proc. IEEE},
63(9), 1975.
\bibitem{greshake2023} K. Greshake et al., ``Not what you've signed up
for: Compromising real-world LLM-integrated applications with indirect
prompt injection,'' \emph{AISec}, 2023.
\bibitem{camel2025} E. Debenedetti et al., ``Defeating prompt
injections by design,'' arXiv:2503.18813, 2025.
\bibitem{thomas2015} P. S. Thomas, G. Theocharous, and M. Ghavamzadeh,
``High confidence policy improvement,'' \emph{ICML}, 2015.
\bibitem{laroche2019} R. Laroche, P. Trichelair, and R. Tachet des
Combes, ``Safe policy improvement with baseline bootstrapping,''
\emph{ICML}, 2019.
\bibitem{geifman2017} Y. Geifman and R. El-Yaniv, ``Selective
classification for deep neural networks,'' \emph{NeurIPS}, 2017.
\bibitem{bates2021} S. Bates, A. Angelopoulos, L. Lei, J. Malik, and
M. I. Jordan, ``Distribution-free, risk-controlling prediction sets,''
\emph{J. ACM}, 68(6), 2021.
\bibitem{angelopoulos2021} A. N. Angelopoulos, S. Bates, E. J.
Cand\`es, M. I. Jordan, and L. Lei, ``Learn then test: Calibrating
predictive algorithms to achieve risk control,'' arXiv:2110.01052, 2021.
\bibitem{leelee2025} K.-H. Lee and S. Lee, ``Machines learn number
fields, but how? The case of Galois groups,'' arXiv:2508.06670, 2025.
\bibitem{anthropic2026} Anthropic Frontier Red Team, ``Patterns and
problems in emerging multiagent systems,'' Aug. 2026.
\bibitem{maurerpontil2009} A. Maurer and M. Pontil, ``Empirical
Bernstein bounds and sample variance penalization,'' \emph{COLT}, 2009.
\bibitem{nds2} NDS2 Lab, Agentic AI Security Initiative, project
index, 2026. \url{https://agentaisec.github.io}.
\end{thebibliography}

\end{document}
